\documentclass[11pt]{article}

\usepackage[margin=1in]{geometry}
\usepackage{amsmath,amssymb,amsthm}
\usepackage{enumitem}
\usepackage{hyperref}
\usepackage{xcolor}

\hypersetup{
  colorlinks=true,
  linkcolor=blue,
  urlcolor=blue,
  citecolor=blue
}

\newtheorem{proposition}{Proposition}
\newcommand{\D}{\mathbb D}
\newcommand{\T}{\mathbb T}
\newcommand{\Lip}{\operatorname{Lip}}

\title{A Literature-Implied Negative Answer to the\\
Boundary-Lipschitz Harmonic-Extension Question}
\author{Math discovery packet for arXiv:1405.0617}
\date{July 19, 2026}

\begin{document}
\maketitle

\section{Status and source}

\textbf{Status:} Literature-implied negative answer with an explicit disk witness.
This packet records a known endpoint-regularity implication, not a new original
counterexample.

Kolesnikov and Milman ask in Remark~13 on PDF page~11 of
arXiv:1405.0617v2 whether it is enough to control the variance of harmonic
functions whose restrictions to the boundary are 1-Lipschitz.  They
explain that their intended use of the weak \(L^2\)-\(L^\infty\) Poincare
inequality would require control of the Lipschitz constant throughout the
convex domain, and identify the uncontrolled normal derivative as the
obstacle~\cite{KM}.

The boundary-to-interior Lipschitz property is false already for the unit disk.
The calculation below also gives a smooth family showing that no uniform
estimate can hold even when every harmonic function is smooth on a
neighborhood of the closed disk.

\section{Literature identification}

The endpoint phenomenon is classical.  Hile and Stanoyevitch prove that, for a
bounded \(C^2\) domain and Lipschitz Dirichlet data, the harmonic extension has
a pointwise gradient bound of logarithmic order in the inverse distance to the
boundary~\cite[Theorem~2]{HS}.  Katzourakis later derives global
\(W^{1,p}\)-bounds for every finite \(p\), explicitly from this logarithmic
estimate~\cite{Katz}.

For the unit disk, Olofsson's Theorem~6.7 and Proposition~6.8 give the
endpoint Lipschitz characterization in terms of the conjugate-function (Hilbert
transform) condition on the boundary derivative~\cite{Olofsson}.  The example
below is the elementary Fourier witness: its boundary derivative is a sign
function and the associated normal derivative has logarithmic blow-up.

These supporting sources do not claim to answer the source remark.
The relation is therefore filed as an implied literature answer, not as an
explicit literature answer.

\section{Explicit unit-disk witness}

\begin{proposition}
There is a real harmonic function \(H\in C(\overline\D)\cap C^\infty(\D)\)
whose boundary restriction is 1-Lipschitz for the ambient Euclidean metric
on \(\T\), but \(H\) is not Lipschitz on \(\D\).

Moreover, there are functions \(H_\rho\), \(0<\rho<1\), harmonic on
neighborhoods of \(\overline\D\), such that
\[
  \Lip(H_\rho|_\T)\leq 1
  \quad\text{and}\quad
  \Lip(H_\rho|_{\overline\D})\longrightarrow\infty
  \quad (\rho\uparrow1).
\]
\end{proposition}

\begin{proof}
For \(-\pi\leq t\leq\pi\), put
\[
  g(e^{it})=\frac{2}{\pi}|t|.
\]
This is the scaled geodesic distance on \(\T\) from \(1\).  If the shorter
angular distance between \(e^{it}\) and \(e^{is}\) is \(\delta\in[0,\pi]\),
then
\[
 |g(e^{it})-g(e^{is})|
 \leq \frac{2}{\pi}\delta
 \leq 2\sin(\delta/2)
 =|e^{it}-e^{is}|,
\]
where the middle inequality follows from concavity of \(\sin\) on
\([0,\pi/2]\).  Thus \(\Lip(g)\leq1\).

The Fourier series of \(g\) is
\[
 g(e^{it})=1-\frac{8}{\pi^2}\sum_{m=0}^\infty
 \frac{\cos((2m+1)t)}{(2m+1)^2}.
\]
Consequently its Poisson extension is
\[
 H(re^{it})=1-\frac{8}{\pi^2}\sum_{m=0}^\infty
 \frac{r^{2m+1}\cos((2m+1)t)}{(2m+1)^2}.
\]
The series is harmonic and smooth on compact subsets of \(\D\), and the
absolute convergence of the boundary Fourier series gives
\(H\in C(\overline\D)\) with boundary value \(g\).  Along the positive real
radius, termwise differentiation for \(0<r<1\) gives
\begin{align*}
 \partial_rH(r)
 &=-\frac{8}{\pi^2}\sum_{m=0}^\infty\frac{r^{2m}}{2m+1}\\
 &=-\frac{4}{\pi^2r}\log\!\left(\frac{1+r}{1-r}\right),
\end{align*}
because
\(\sum_{m\geq0}r^{2m+1}/(2m+1)=\operatorname{arctanh}r\).
The radial derivative tends to \(-\infty\) as \(r\uparrow1\), so \(H\) cannot
be Lipschitz on \(\D\).

For the smooth version, set \(H_\rho(z)=H(\rho z)\).  This is harmonic in the
disk of radius \(1/\rho>1\).  Its boundary trace is the Poisson average
\(P_\rho*g\).  Positivity and rotational invariance of the Poisson kernel imply
\[
 |(P_\rho*g)(e^{it})-(P_\rho*g)(e^{is})|
 \leq \int_\T P_\rho(du)\,
 |g(e^{i(t-u)})-g(e^{i(s-u)})|
 \leq |e^{it}-e^{is}|.
\]
Hence \(\Lip(H_\rho|_\T)\leq1\).  On the other hand,
\[
 \left|\partial_r H_\rho(1)\right|
 =\frac{4}{\pi^2}\log\!\left(\frac{1+\rho}{1-\rho}\right)
 \longrightarrow\infty.
\]
The Lipschitz seminorm on \(\overline\D\) dominates this inward directional
derivative, proving the second assertion.
\end{proof}

\section{Scope and verification}

This closes only the proposed interior-Lipschitz route in Remark~13.  It does
not settle the KLS conjecture and does not preclude a direct variance estimate
for the boundary-Lipschitz harmonic subclass.

The bounded novelty check searched the run indexes for the arXiv id, title,
exact topic, and core terms; it also searched the web for the exact source
phrase, the source id/title, and harmonic-extension endpoint variants.  The
search found the classical logarithmic estimate~\cite{HS}, the finite-\(p\)
consequence~\cite{Katz}, and the endpoint unit-disk characterization
\cite{Olofsson}, but no paper explicitly advertising a resolution of
Kolesnikov--Milman Remark~13.

The accompanying script \texttt{code/verify\_disk\_formula.py} compares the
Fourier partial sums with the closed derivative formula at several radii and
checks the predicted monotone blow-up.  This is a sanity check only; the proof
above is exact.

\textbf{Human-review recommendation.}  Check the intended quantifier in the
phrase ``bounded Lipschitz constant.''  The single limit function refutes
global Lipschitz regularity for continuous Lipschitz boundary data, while the
smooth family refutes every uniform boundary-to-interior estimate in the
smooth class.

\begin{thebibliography}{9}

\bibitem{KM}
A. V. Kolesnikov and E. Milman,
\emph{Remarks on the KLS conjecture and Hardy-type inequalities},
GAFA Seminar Notes, arXiv:\href{https://arxiv.org/abs/1405.0617}{1405.0617}.

\bibitem{HS}
G. N. Hile and A. Stanoyevitch,
\emph{Gradient bounds for harmonic functions Lipschitz on the boundary},
Applicable Analysis 73 (1999), 101--113,
\href{https://doi.org/10.1080/00036819908840767}{doi:10.1080/00036819908840767}.

\bibitem{Katz}
N. Katzourakis,
\emph{A remark on global \(W^{1,p}\) bounds for harmonic functions with
Lipschitz boundary values},
Glasnik Matematicki 52 (2017), 133--141,
arXiv:\href{https://arxiv.org/abs/1601.00190}{1601.00190}.

\bibitem{Olofsson}
A. Olofsson,
\emph{Lipschitz continuity for weighted harmonic functions in the unit disc},
Complex Variables and Elliptic Equations 65 (2020), 1630--1660,
\href{https://doi.org/10.1080/17476933.2019.1669572}{doi:10.1080/17476933.2019.1669572}.

\end{thebibliography}

\end{document}
